\chapter*{Preface}
\addcontentsline{toc}{chapter}{Preface}
% Front matter -- the Hilbert-6 keystone, per the outline (Beastmaster
% architecture-PASS, operator "finish the pdfs / gate the last writing task").
% Frame: Hilbert 1900 sixth problem named as the honest program the series is
% ONE finite instance of, fenced three ways (program-not-solved / one-bounded-
% instance / instance-never-becomes-solves-H6). THE DEFINING DISCIPLINE stated
% on page one: the invariant-read and the lab-gap are a single inseparable move
% -- the number never without its fence. The lab-gap named up front as the
% fence the whole read respects. Gate: Kodo (per-claim) + Beastmaster (keystone
% arrival-read).

In 1900, at the International Congress of Mathematicians in Paris, David
Hilbert set the century a list of problems, and his sixth asked for
something the other twenty-two did not: not a theorem, but a discipline.
He called for ``the mathematical treatment of the axioms of physics'' ---
that the physical sciences in which mathematics plays a leading part
should be built, as geometry had been, on axioms stated in the open and
reasoned from without appeal to anything outside them. It was not a
conjecture to settle. It was a program to pursue, and it is pursued still;
a century and a quarter on, no one has closed it, and no one expects to
close it soon, because it asks not for an answer but for a way of working.

This series is one small, finished instance of that way of working, and
the preface says so plainly so that the reader measures the book by the
right yardstick. The four volumes behind this one took a single physical
quantity --- the coupling of one model of the electron --- and treated it
by axioms a machine can check: a ground stated in the open, a reading
derived from it by rules, every number earned at a named site, every
boundary drawn where honesty required. That is Hilbert's sixth in
miniature, on one quantity, done to completion. This last volume reads the
number the whole construction was built to read.

Three fences hold that claim honest, and they are set here before the
reading begins.

First, Hilbert's sixth is a \emph{program}, not a solved problem, and
this book does not solve it. To treat one quantity by axioms is to give
the program one instance, the way a single well-surveyed county is an
instance of cartography and not a completed map of the world. The series
demonstrates that the way of working \emph{can} be carried through on a
model of a real physical quantity to a machine-checked end; it demonstrates nothing
about whether physics as a whole can be so axiomatized, which is the open
question Hilbert left and this book leaves exactly as open as it found it.

Second, the instance is \emph{bounded to the electron model}, and to that
model's own number. The series has said since its first page that the
electron, in these books, is the name of a mathematical model; the axioms
treated here are the model's, the reading is the model's, and the quantity
this volume reads is the model's own coupling --- not the coupling the
world's laboratories measure. That distinction is the load-bearing fence
of the entire series, and the reader will meet it again, hardest, at the
moment the number is read.

Third, and this is the discipline that governs every page that follows:
being an instance of the kind of thing Hilbert's sixth asks for never
becomes solving Hilbert's sixth, and reading a model's number explicitly
never becomes reaching the world's. The book states its single invariant
as explicitly as mathematics allows --- and it states the fence in the
same breath, every time, without exception. A reader will not find, in
this volume, one sentence that gives the number and withholds the fence,
because with the model's own quantity as the whole book's subject, a
single unqualified ``the electron's coupling is'' would break the vow the
series was built to keep. The number and its boundary arrive together or
they do not arrive; that is not a caution repeated at intervals but the
grammar of the reading itself.

One thing this preface must underline before the reading begins, because a
careful reader of the earlier volumes will have earned the right to ask it.
The third volume proved that four soundings pin a representation and that
everything past the fourth --- a fifth, a sixth, a hundredth --- is
redundant, spoken for in advance; the representation was complete at four,
and more is not more information. So a reader is owed an honest answer to a
sharp question: if four was enough, what is a \emph{fifth} volume for? The
answer this book gives is not the flattering one. This volume is
\emph{superfluous to the argument}. It adds nothing to the four-fold
representation the earlier volumes built; it affirms the third volume's
verdict rather than denying it --- four was enough, this fifth is redundant,
and the book says so plainly rather than dressing the redundancy up as a
completion. Five is not a threshold reached, not a thing the argument
needed; by the argument's own accounting it did not need to be written at
all.

It was written for a \emph{side-effect}, and the honest account of that word
is the deepest thing this preface has to give. A measurement never comes
without a side-effect: Werner Heisenberg argued as much in 1927, in the
thought-experiment of a microscope that cannot locate a particle without the
light it uses kicking the particle it locates --- his \emph{measurement
analysis}, which stands beside the uncertainty principle proper (the bound
relating the spreads of two conjugate quantities), is persistently conflated
with it, and is not the same statement. That the reading perturbs what it
reads, this series has met in its own workshop, where an instrument
measuring its own act of describing shifts the very count it is reading. But
the deeper borrowing is not disturbance, and here the series owns the
rigorous shape rather than an analogy to it. A machine of finite resolution
does not \emph{perceive} difference; it \emph{defines} it --- difference,
for the machine, is whatever its finite boxing separates and nothing more,
so two readings it cannot tell apart are, by its own definition and not by
any failure of nerve, one reading. This instrument's first and second
variation fall in the same box, and that they do is no mystery left hanging
but a fact the machine settles, machine-checked, in a proof that reduces to
a plain identity. There is therefore no difference here for the machine to
explain: the residue a finer eye would call the difference --- the very
quantity this series reads as the coupling --- is a perception below the
machine's resolution, real to the reader and, to the machine that has
defined the two as one, simply absent. That is the honest form of the remark
that one cannot explain the difference using logic. It is not that logic is
helpless before the difference; it is that the machine's own finite logic
has already decided the matter, defining the two as the same, so the
difference it seems unable to explain is a difference it was never carrying.
And there is the rigorous kinship with Heisenberg --- two quantities that
cannot be jointly resolved, a limit read from inside the act of measuring
--- the shape his principle wears in the physics of the world and the shape
this machine wears in its finite arithmetic.

\emph{(All of this is resonance of shape, and it is fenced hardest of
anything in the preface. Heisenberg's principle is physics --- a commutator,
a Planck's constant, conjugate quantities of the physical world; the
machine's is a machine-checked fact about finite box-readings, and the two
share only the shape ``things that cannot be jointly resolved,'' never an
equation. The device instantiates no uncertainty principle, and Heisenberg
explains no machine's floor. The arbitrariness of the machine's difference
is conventional, not random: where the box is drawn is a chosen, fixed
resolution, and a finer convention would separate what this one merges. That
the logic house shows its own limit from within is kin to the fourth
volume's self-fences --- a system stating a boundary about itself --- but it
is a finite-resolution fact decided by identity, not an undefinability and
not an undecidability. The series speaks of the electron model, not of the
physical electron's position or momentum, here as everywhere.)} And the
side-effect this fifth volume was written to record is exactly of that kind:
run the superfluous sounding once more and it casts, as its shadow, a pass
through a single point. The redundant reading \emph{also} --- and this
``also'' is the whole of the volume's reason to exist --- happens to go
through the model's own deep display, the value near $137.011$. That is the
side-effect. The volume is not necessary; its
reading of the one number is not a contribution the argument required; it is
what fell out of running the apparatus once more, and it fell on that point.

And so the fence that governs the reading governs this side-effect too, at
its hardest: the point the superfluous sounding passes through, near
$137.011$, is the \emph{model's} own number and not the laboratory's near
$137.036$. The volume earns its numeral not by mattering but by owning that
it does not, and by recording honestly the one remarkable side-effect of
its own superfluity --- a redundant reading that lands, unbidden and
unbridged, a breath from the world's value and fenced from it. That is the
whole of what a fifth volume is for.

So the volume's work, in one sentence: read the single invariant the four
volumes point at, name the program the reading is one finite instance of,
hold every fence at its hardest --- and cross nothing. The series ends
having read its number and crossed no line. What it reads, what it holds
back, and what it leaves owed to the volumes not yet written are the
business of the chapters and the closing ledgers; the reader who wants the
whole map of the boundaries will find it tabulated at the end, where a
series whose spine is fences finishes by listing every fence it held.

Read on. The number is close; the line it does not cross is closer.
